% ============================================================================
% USB MODES, RE-ENUMERATION, AND BENCH JTAG INTERACTION
% ============================================================================

\section{USB Modes and Bench Behavior}

\subsection{Transport Selection at the Gateway}

The Go gateway picks a transport via the \texttt{TransportMode} enum
defined in \texttt{star-gateway/internal/manager/types.go}. Four modes
are valid:

\begin{center}
\begin{tabular}{|l|l|l|}
\hline
\textbf{Mode}              & \textbf{Wire-format invariant}        & \textbf{When to use}                                       \\
\hline
\texttt{auto}              & Same framed protocol, any transport   & Default; prefers USB CDC, falls back to SPI                \\
\texttt{force-usb}         & Full HARQ + reset handshake on USB    & Strict-host environments where SPI is intentionally off    \\
\texttt{simple-usb}        & Framed protocol, no Layer-3 HARQ      & Pi 5 xHCI hosts where USB CDC is already lossless          \\
\texttt{force-spi}         & Full HARQ over SPI; USB hot-plug off  & Closed-loop motor-control rigs that need bounded jitter    \\
\hline
\end{tabular}
\end{center}

\subsubsection{Simple USB mode}

\textbf{What it is:} The gateway opens the USB CDC port (typically
\texttt{/dev/ttyACM\{1,2,3\}} on PROD or \texttt{/dev/ttyACM\{4,5,6\}}
on TOM) and treats the byte stream as already-reliable. It performs:

\begin{itemize}
  \item CRC-32 frame validation as defense in depth.
  \item No reset handshake at the link layer.
  \item No strict sequence-validation (gaps are tolerated rather than
        triggering NACK + retransmit storms).
  \item Relaxed health monitoring (lower retry budgets, longer dead
        timers).
\end{itemize}

\textbf{Why it is safe:} USB 2.0 Full-Speed already delivers a CRC-16
on every packet and re-transmits transparently in the host controller.
Layering HARQ on top is wasted work and introduces latency that is
nontrivial under typical CDC ACM scheduling (every-frame ACK + retry
budget). On a Pi 5 xHCI host, simple-usb has been bench-verified to
enumerate cleanly and survive arbitrary disconnect/reconnect cycles.

\textbf{When to avoid:} Strict-host environments (some Windows HCDs,
some macOS releases, embedded controller hubs) may be picky enough
about USB framing or clock accuracy that the loss-tolerant assumption
fails. In that case prefer \texttt{force-usb} (or fall back to SPI).

\subsubsection{Complex (full) USB mode -- \texttt{force-usb}}

This is the original transport design: a full link-layer state machine
with reset handshake (RST/RST-ACK), strict 16-bit sequence checking,
NACK-driven retransmits with Chase-Combining HARQ, and deadline
monitors. It is the canonical mode on hardware that also has the SPI
link active and where deterministic recovery is more important than
raw throughput.

\subsubsection{Forcing SPI -- \texttt{force-spi}}

Used by the ROS2 motor-control bridge (\texttt{star\_spi\_bridge}).
USB hot-plug is disabled so that wandering USB cables can never knock
the closed-loop control off the deterministic SPI path. The full HARQ
stack runs on this transport because SPI itself has no error
correction and PWM/motor noise is real.

\subsection{USB enumeration topology}

The RX72N exposes a single composite USB device that the host sees as
three CDC-ACM serial endpoints in a single \texttt{idVendor / idProduct}
tuple. Bench-verified mappings (commit \texttt{341bd09ce} on
\texttt{feat/usb0}, 2026-04-24):

\begin{itemize}
  \item \textbf{PROD board} (E2L \texttt{5AS071312B}, 24~MHz crystal):
        \texttt{/dev/ttyACM1} = PROTO,
        \texttt{/dev/ttyACM2} = DECODED,
        \texttt{/dev/ttyACM3} = LOG.
  \item \textbf{TOM board} (E2L \texttt{5AS071321B}, HOCO PLL):
        \texttt{/dev/ttyACM4} = PROTO,
        \texttt{/dev/ttyACM5} = DECODED,
        \texttt{/dev/ttyACM6} = LOG.
  \item Both enumerate as \texttt{idVendor=045b idProduct=0235}, and
        both heartbeat on all three ports under \texttt{p0/p1/p2}.
\end{itemize}

\subsection{Re-enumeration: when ttyACM* changes}

The OS assigns \texttt{/dev/ttyACM*} indices in order of appearance.
Anything that drops the device and re-enumerates it can shift the
mapping. Common causes on the bench:

\begin{itemize}
  \item Reflash via E2 Lite -- the firmware re-init cycles USB power
        and the USBb peripheral, so the host sees a disconnect followed
        by a fresh enumerate. The new \texttt{ACM} numbers usually
        come back in the same order, but if another USB-CDC device
        was plugged in between cycles, indices can advance.
  \item Power-cycling the board (USB unplug or barrel-jack toggle).
  \item Hot-plugging another CDC ACM device (an Arduino, an ESP32,
        another RX72N) while the gateway is running.
\end{itemize}

To make the assignment stable, install
\texttt{scripts/71-star-symlinks.rules} via
\texttt{scripts/setup-pi5-host.sh}. That udev rule creates persistent
symlinks (e.g.\ \texttt{/dev/star-rx72n-proto},
\texttt{/dev/star-rx72n-decoded}, \texttt{/dev/star-rx72n-log}) keyed
on the device's serial / interface index, so the gateway can open
those symlinks regardless of what \texttt{/dev/ttyACM*} number the
kernel happens to pick on a given boot.

\subsection{Bench JTAG interaction (E2 Lite vs CY7C65213)}

On the STAR bench, the Renesas E2 Lite debugger and the Cypress
CY7C65213 USB-UART bridge are mutually exclusive on the host:

\begin{itemize}
  \item While the E2 Lite is attached and holding the JTAG/FINE
        interface, the CY7C65213 does \emph{not} enumerate as
        \texttt{/dev/cu.usbmodem*} or
        \texttt{/dev/cu.usbserial-*}. This is bench-harness wiring
        (shared USB power gating), not a fault.
  \item After flashing, physically disconnect the E2 Lite to expose
        the bridge to the host before running picocom or starting
        the gateway against the SCI9 fallback path.
  \item The USB0 path (\texttt{/dev/ttyACM*}) is not affected by E2
        Lite presence -- it runs on the RX72N USBb peripheral on
        the dedicated USB-C connector and is independent of SCI9.
\end{itemize}

\subsection{Clock-source notes (TOM vs PROD)}

The TOM board has no external main crystal and runs USB0 from
HOCO~$\rightarrow$~PLL (PLLSRCSEL=1, multiplier 12,
\texttt{k\_pll\_multiplier\_12\_hoco = 0x1710}). The Renesas HUM
explicitly recommends an 8--24~MHz crystal as the USB clock source
(HUM page~345 PLLCR Note~1, section~9.1, Table~63.18). HOCO is
$\pm$2.44~\% accurate at $T_a \ge -20^\circ$C; USB 2.0 Full-Speed is
specified at $\pm$0.25~\%, so HOCO is roughly 10x off-spec on paper.

In practice, the Pi~5's xHCI controller accepts the HOCO-driven clock
without complaint, and TOM enumerates the 3-port CDC composite
cleanly. Treat the HUM warnings as a portability risk (some macOS
revisions or Windows HCDs may reject the clock source) rather than a
blocking defect for ship-to-Pi deployments. For field deployments on
unknown hosts, prefer the 24~MHz crystal-clocked PROD configuration.

\subsection{References}

\begin{itemize}
  \item \texttt{star-gateway/internal/manager/types.go} -- transport
        mode definitions.
  \item \texttt{star-rx72n-firmware/libs/rx\_usb/} -- USBb peripheral
        driver.
  \item \texttt{scripts/71-star-symlinks.rules} -- udev rules for
        stable \texttt{/dev/star-rx72n-*} symlinks.
  \item \texttt{scripts/setup-pi5-host.sh} -- one-time Pi 5 host
        setup that installs the rules above plus the USB
        autosuspend overrides.
\end{itemize}
